% !PRINTED 133
\weberpage{133}

Wenn die absoluten Werthe von $\varepsilon_1,\varepsilon_2,
\ldots$ unter hinlänglich kleinen Werthen liegen, so können wir von der Function $f_1(x)$ zunächst beweisen,
\begin{center}
\gesperrt{dass sie keine Wurzeln ausserhalb der Gebiete $(\varrho),(\varrho'),(\varrho''),\ldots$ hat,}
\end{center}
und zweitens,
\begin{center}
\gesperrt{dass die Anzahl der Wurzeln von $f_1(x)$ innerhalb $(\varrho)$ genau $a$, innerhalb $(\varrho')$ genau $b$, innerhalb $(\varrho'')$ genau $c$ u. s. f. beträgt.}
\end{center}
Bei dem letzten Theil des Satzes ist aber zu beachten, dass, wenn $f_1(x)$ mehrfache Wurzeln hat, diese nach ihrer Vielfachheit gezählt werden müssen.

Wir construiren nach \S 31, IV in der Ebene $x$ einen Kreis mit dem Radius $R$ und dem Nullpunkt als Mittelpunkt, der die Gebiete $(\varrho),(\varrho'),(\varrho''),\ldots$ einschliesst, so dass ausserhalb dieses Kreises keine Wurzeln von $f_1(x)$ mehr liegen, und bezeichnen das innerhalb dieses Kreises, aber ausserhalb $(\varrho),(\varrho'),(\varrho''),\ldots$ liegende Gebiet mit $G$. Es ist dazu noch zu bemerken, dass $R$ von den $\varepsilon_1,\varepsilon_2,
\ldots$ unabhängig angenommen werden kann, so lange für die absoluten Werthe dieser Grössen eine bestimmte obere Grenze festgesetzt wird.

% Exact unretouched crop from the registered Volume-I scan, printed Fig. 4.
\weberfigureimage{weber-b1-fig04.png}{0.42}

Ist nun $x$ ein Punkt des Gebietes $G$, so ist der absolute Werth von $x-\alpha$ grösser als $\varrho$, und mithin nach (2)
\begin{equation}
 |f(x)|>\varrho^a\varrho'^b\varrho''^c\cdots .
\tag{5}
\end{equation}
Ist nun $\alpha_1$ eine Wurzel von $f_1(x)$, so folgt aus (4)
\begin{equation}
 f(\alpha_1)=-\varphi(\alpha_1),
\tag{6}
\end{equation}
und $\alpha_1$ kann daher nicht in dem Gebiete $G$ liegen, wenn man dafür sorgt, dass innerhalb $G$ überall
\begin{equation}
 |\varphi(x)|<\varrho^a\varrho'^b\varrho''^c\cdots,
\tag{7}
\end{equation}
was durch genügende Verkleinerung der oberen Grenze von $\varepsilon_1,\varepsilon_2,
\ldots$ immer möglich ist. Hiermit ist der erste Theil unserer Behauptung erwiesen, dass, wenn die Ungleichung (7) befriedigt ist, keine Wurzeln von $f_1(x)$ ausserhalb der Gebiete $(\varrho),(\varrho'),(\varrho'')\ldots$ liegen.
